\section*{الفصل \ref{ch:b1:clastic-sediments} --- الرسوبيات الفتاتية وبيئاتها}

\begin{solution}{exo:b1:clastic-sediments:1}
$\varphi = -\log_2 d$: \qty{0.18}{mm} تعطي $2.47$، رملًا ناعمًا (2 إلى 3)؛ \qty{0.75}{mm} تعطي $0.42$، رملًا خشنًا (0 إلى 1)؛ \qty{3}{mm} تعطي $-1.58$، حصاة ($-1$ إلى $-2$)؛ \qty{0.012}{mm} تعطي $6.38$، طميًا ناعمًا (6 إلى 7)؛ \qty{150}{mm} تعطي $-7.23$، حصًى كبيرًا ($-6$ إلى $-8$).
\end{solution}

\begin{solution}{exo:b1:clastic-sediments:2}
$M_z = (1.05 + 1.71 + 2.43)/3 = 1.73\,\varphi$، أي $2^{-1.73} =
\qty{0.30}{mm}$، رمل متوسط. $\sigma_I = (2.43 - 1.05)/4 + (2.98 - 0.62)/6.6
= 0.345 + 0.358 = 0.70\,\varphi$: فرز جيّد نسبيًّا (0.50 إلى 0.71)، دون حدّ \omterm{def:b1:clastic-sediments:sorting}{الفرز} المتوسط بقليل.
\end{solution}

\begin{solution}{exo:b1:clastic-sediments:3}
$q_b = (1 - 0.4) \times 0.02 \times (\num{0.01}/60)/2 = \qty{1.0e-6}{m^2.s^{-1}}$
من الحبيبات الصلبة. في اليوم، $\num{1.0e-6} \times 86400 = \qty{0.0864}{m^3}$ لكل متر من القمّة،
أو $0.0864 \times 2650 = \qty{229}{kg}$ من الكوارتز.
\end{solution}

\begin{solution}{exo:b1:clastic-sediments:4}
(أ)~50 عقدة تساوي $50 \times 1852/3600 = \qty{25.7}{m.s^{-1}}$؛ و\num{12} عقدة،
\qty{6.17}{m.s^{-1}}. (ب)~$\qty{100000}{m}/\qty{6.17}{m.s^{-1}} =
\qty{16200}{s}$، نحو \qty{4.5}{h}. (ج)~480 ميلًا تساوي \qty{772}{km}، و\num{400} ميل
\qty{644}{km}.
\end{solution}

\begin{solution}{exo:b1:clastic-sediments:5}
(أ)~$(\varphi_{84} - \varphi_{16})/2 = 2 \times 0.9945\,\sigma/2 =
0.9945\,\sigma$، و$(\varphi_{95} - \varphi_5)/3.3 = 2 \times
1.6449\,\sigma/3.3 = 0.9969\,\sigma$: كلاهما يقدّر $\sigma$ في حدود ثلث في المئة. (ب)~لا يلزم أن يكون الرسوب الحقيقي طبيعيًّا في وحدات فاي: فخليط صغير من حبيبات أخشن بكثير أو أدقّ بكثير يغيّر $\varphi_5$ أو $\varphi_{95}$ كثيرًا، و$\varphi_{16}$ و$\varphi_{84}$ قليلًا، فيكبر التقدير الثاني بينما لا يكبر الأول. والاثنان يختلفان أكثر ما يختلفان في الخلائط، كالرمل مع بعض الحصى والطين، أو الرسوب ثنائي المنوال.
\end{solution}

\begin{solution}{exo:b1:clastic-sediments:6}
(أ)~$\sin\theta = T/L = 1.5/15 = 0.1$، إذن $\theta = \ang{5.74}$. (ب)~$a =
c\tan\theta = 0.5 \times 0.1005 = \qty{0.050}{cm}$ في الدقيقة، أو
\qty{3.0}{cm} في الساعة. (ج)~طبقة تُرفع بسنتيمترات في الساعة تتلقّى رملًا يهبط من التعليق بسرعة كبيرة جدًّا: تيار محمّل بكثافة يتباطأ، كفيضان متراجع أو ذيل تيار عكارة.
\end{solution}

\begin{solution}{exo:b1:clastic-sediments:7}
(أ)~$(350 + 20 + 5 + 340 + 15 + 30)/6 = 760/6 = 126.7$: اتجاه جنوبي شرقي، مع أن كل قراءة تقع ضمن \ang{20} من الشمال. (ب)~$S =
\sum\sin\alpha_i = 0.672$، $C = \sum\cos\alpha_i = 5.692$، إذن $\bar\alpha =
\operatorname{atan2}(0.672, 5.692) = \ang{6.7}$، و$\bar R =
\sqrt{0.672^2 + 5.692^2}/6 = 5.732/6 = 0.955$. (ج)~المتوسط الحسابي يعامل 350 و\num{005} كأنهما متباعدان، بينما هما على الدائرة متباعدان ب\ang{15} فقط: يجب أن تُحسب السموت كمتجهات وحدوية. و$\bar R$ الكبير يدلّ على أن القراءات متجمّعة بإحكام حول \ang{7}.
\end{solution}

\begin{solution}{exo:b1:clastic-sediments:8}
الحجر A فيه أكثر من \qty{5}{\percent} من الطين: غير ناضج، وغير ناضج تركيبيًّا أيضًا، مع \omterm{def:b1:rock-forming-minerals:groups}{الفلسبار} وأجزاء صخرية؛ ترسّب بسرعة قرب مصدره، على مروحة فيضية أو سهل فيضي، أو بتيار عكارة. والحجر B فيه أقلّ من \qty{5}{\percent} من الطين، وفرزه جيّد ($\sigma_I < 0.5\,\varphi$)، لكن حبيباته شبه زاوية ($\rho < 3$): ناضج، كقضيب نهري أو رمل رفّ غربلته التيارات. أما الحجر C فلا طين فيه، وفرزه جيّد جدًّا، وحبيباته مستديرة ($\rho > 3$): فائق النضج، ومع \qty{98}{\percent} كوارتزًا فهو ناضج تركيبيًّا أيضًا، كما على شاطئ أو في حقل كثبان صحراوية.
\end{solution}

\begin{solution}{exo:b1:clastic-sediments:9}
(أ)~$1.5/\sqrt{9.81 \times 0.3} = 0.87$؛ $1.0/\sqrt{9.81 \times 2} = 0.23$.
(ب)~$h = U^2/g = 1.5^2/9.81 = \qty{0.23}{m}$؛ $U = \sqrt{9.81 \times 2} =
\qty{4.43}{m.s^{-1}}$. (ج)~يكبر عدد فرود كلما قلّ العمق عند سرعة معيّنة: الصفائح الرقيقة السريعة من الماء، مثل تراجع موجة على شاطئ، تبلغ بسهولة أعدادًا قرب 1، أما نهر عمقه عدّة أمتار فيحتاج إلى عدّة أمتار في الثانية، وهذه سرعة لا يبلغها معظم الأنهار إلا في الفيضانات.
\end{solution}

\begin{solution}{exo:b1:clastic-sediments:10}
(أ)~مع النقاط الوسطى $-0.5, 0.5, \ldots, 4.5$: $m = (-1 + 4 + 45 + 100 + 52.5
+ 22.5)/100 = 2.23\,\varphi$؛ مجموع المربعات المنحرفة الموزونة بالنسب المئوية هو \num{107.71}، ف$s = \sqrt{1.0771} = 1.04\,\varphi$:
فرز ضعيف. (ب)~النسبة المئوية الأخشن تراكميًّا هي 0 و\num{2} و\num{10} و\num{40} و\num{80} و\num{95} و\num{100} عند $-1, 0, \ldots, 5\,\varphi$. إذن $\varphi_5 = 0 + 3/8 = 0.375$،
$\varphi_{16} = 1 + 6/30 = 1.20$، $\varphi_{50} = 2 + 10/40 = 2.25$،
$\varphi_{84} = 3 + 4/15 = 3.27$، $\varphi_{95} = 4.0$. $M_z = (1.20 + 2.25 +
3.27)/3 = 2.24\,\varphi$ و$\sigma_I = (3.27 - 1.20)/4 + (4.0 - 0.375)/6.6 =
0.517 + 0.549 = 1.07\,\varphi$. (ج)~تتفق الطريقتان بفارق نحو \num{0.01} في المتوسط و\num{0.03} في \omterm{def:b1:clastic-sediments:sorting}{الفرز}: رمل ناعم ($2.24\,\varphi$ تساوي \qty{0.21}{mm})، فرزه ضعيف.
\end{solution}

\begin{solution}{exo:b1:clastic-sediments:11}
(أ)~$U = \sqrt{0.02 \times 9.81 \times 50 \times 0.005/0.003} =
\qty{4.04}{m.s^{-1}}$، أو 7.9 عقدة. (ب)~$(\Delta\rho/\rho)\,h = U^2
C_f/(gS) = 25.7^2 \times 0.003/(9.81 \times 0.005) = \qty{40}{m}$: مثلًا تيار سمكه \qty{400}{m} مع فائض كثافة \qty{10}{\percent}، وهذا يبيّن مدى ضخامة التيار وكثافته في 1929. (ج)~تتناسب السرعة مع $\sqrt S$: $4.04/\sqrt 5 = \qty{1.81}{m.s^{-1}}$. التيار الأبطأ تكون اضطراباته أضعف، فيُسقط حبيباته الأخشن؛ ويصير أقلّ كثافةً، ويتباطأ أكثر، ويفرد حمولته على السهل طبقةً متدرّجة.
\end{solution}

\begin{solution}{exo:b1:clastic-sediments:12}
(أ)~$0.7 \times \qty{5e10}{kg}/\qty{2000}{kg.m^{-3}} = \qty{1.75e7}{m^3}$ في
السنة. (ب)~موزّعًا على واجهة ارتفاعها \qty{30}{m} وعرضها \qty{100}{km}، يدفع الواجهة بمقدار $\num{1.75e7}/(30 \times \num{1e5}) = \qty{5.8}{m}$ في
السنة. (ج)~$a = v\tan\alpha = 5.83 \times \tan\ang{0.5} = \qty{0.051}{m}$، نحو
\qty{5}{cm} في السنة، أي أقلّ بمئة مرّة من التقدّم؛ ويتكوّن التعاقب بسمك \qty{30}{m} في $30/0.051 = 590$ سنة.
\end{solution}

\begin{solution}{pb:b1:clastic-sediments:1}
\textbf{1.} الوحدة~1: $(5.0 + 6.6 + 8.4)/3 = 6.67\,\varphi$، \qty{0.0098}{mm}
(\qty{9.8}{\micro m})؛ الوحدة~2: $3.47\,\varphi$، \qty{0.090}{mm}؛ الوحدة~3:
$1.80\,\varphi$، \qty{0.29}{mm}؛ الوحدة~4: $0.73\,\varphi$، \qty{0.60}{mm}.

\textbf{2.} الوحدة~1: $(8.4 - 5.0)/4 + (9.6 - 4.2)/6.6 = 0.850 + 0.818 =
1.67\,\varphi$؛ الوحدة~2: $0.550 + 0.545 = 1.10$؛ الوحدة~3: $0.275 + 0.273 = 0.55$؛
الوحدة~4: $0.725 + 0.712 = 1.44$.

\textbf{3.} الوحدة~1: طمي ناعم، فرز ضعيف؛ الوحدة~2: رمل ناعم جدًّا، فرز ضعيف؛ الوحدة~3: رمل متوسط، فرز جيّد نسبيًّا؛ الوحدة~4: رمل خشن، فرز ضعيف.

\textbf{4.} الوحدة~3. عند مصبّ النهر يحرّك التيار والأمواج الرمل مرارًا ويحملان الدقائق بعيدًا، بينما يترسّب طين ما قبل الدلتا من التعليق خليطًا من كل الأحجام الدقيقة التي تصل إليه.

\textbf{5.} الحجم $8\,\varphi$ يقع على بُعد $1.4/1.8$ من الطريق من
$\varphi_{50}$ إلى $\varphi_{84}$، ف$50 + (1.4/1.8) \times 34 = 76$ في المئة من
العيّنة أخشن منه، ونحو \qty{24}{\percent} بحجم الطين. وتترسّب مثل هذه الجسيمات الدقيقة من الماء الشبه الراكد وحده: فالموقع كان تحت مدى الأمواج والتيارات.

\textbf{6.} تحمل القناة الحصى، المتدحرج على قاعها في الفيضانات، والرمل معًا؛ والرسوب يخلط حصى قاع التيار الأعمق مع رمل التيارات الأضعف، فيغطّي فئات كثيرة.

\textbf{7.} يكبر المتوسط من \qty{0.0098}{mm} إلى \qty{0.090}{mm} و\qty{0.29}{mm}: التعاقب متزايد الخشونة نحو الأعلى.

\textbf{8.} الوحدة~1، طين مصفّح: طين ما قبل الدلتا؛ الوحدة~2، طمي ورمل متناوبان: الواجهة السفلى للدلتا؛ الوحدة~3، رمل متطبّق متقاطع: حواجز المصبّ في الواجهة العليا للدلتا.

\textbf{9.} وصلت كل منطقة إلى الموضع الذي كانت فيه أمامها سابقًا: بنت الدلتا إلى الأمام، نحو البحر، فوق الموقع.

\textbf{10.} \qty{21}{m} على الأقل، وهو سمك الوحدات 1 إلى 3، لأن التعاقب ملأ الماء من أرضية طين ما قبل الدلتا إلى حاجز المصبّ قرب السطح؛ وأكثر من ذلك، لأن الطين انضغط منذ ذلك (\cref{ch:b1:diagenesis-basins}).

\textbf{11.} قناة توزيع من سهل الدلتا، تقطع حاجز المصبّ مع تقدّم الدلتا إلى الأمام؛ وحصاها هو \omterm{def:b1:erosion-transport:load}{الحمولة القاعية}.

\textbf{12.} لا: الحاجز الانعطافي يبدأ بسطح تآكل وحصى بقايا، ويتناقص حجمه نحو الأعلى حتى طين السهل الفيضي. أما واجهة الدلتا فتتزايد خشونتها نحو الأعلى من الطين.

\textbf{13.} $\qty{16}{m}/\num{4000}\text{ سنة} = \qty{4.0}{mm}$ في السنة.

\textbf{14.} $v = a/\tan\alpha = 0.0040/\tan\ang{0.5} = 0.0040/0.008727 =
\qty{0.458}{m}$ في السنة.

\textbf{15.} $0.458 \times 4000 = \qty{1830}{m}$، نحو \qty{1.8}{km}.

\textbf{16.} $\qty{21}{m}/\qty{4.0}{mm.yr^{-1}} = 5250$ سنة.

\textbf{17.} تتلقّى واجهة الدلتا معظم الرمل قرب المصبّات، ويتلقّى طين ما قبل الدلتا أقلّ بكثير، وسهل الدلتا قليلًا جدًّا إلا في الفيضانات؛ والمصبّات النشطة تتحرّك أيضًا من مكان إلى آخر، فقد تتلقّى نقطة كثيرًا لفترة ثم تكاد لا تتلقّى شيئًا. والمعدّل الواحد بين طبقتين مؤرّختين متوسّط عبر هذه التغيّرات، وقد فقد الطين جزءًا من سمكه بالانضغاط.

\textbf{18.} المجموع 1544، فالمتوسط الحسابي \ang{128.7}، أي اتجاه جنوبي شرقي: لا معنى له، لأن القراءات تتشتّت حول الشمال، وتُحسب القيم القريبة من 340 والقريبة من 010 كأنها بعيدة عن بعضها.

\textbf{19.} $S = \sum\sin\alpha_i = 1.698$ و$C = \sum\cos\alpha_i =
11.248$.

\textbf{20.} $\bar R = \sqrt{1.698^2 + 11.248^2}/12 = 11.376/12 = 0.948$: القراءات متجمّعة بإحكام حول اتجاه واحد.

\textbf{21.} جرت الأنهار، وبنت الدلتا، نحو الشمال تقريبًا (\cref{prop:b1:clastic-sediments:azimuths})؛ وكانت اليابسة إلى الجنوب.

\textbf{22.} لأن الواجهات الأمامية قيست على طبقات لم تعد أفقية: فاتجاهات \omterm{def:b1:geological-maps:attitude}{الانحدار} لها هي اتجاهات الطبقات المائلة مضافًا إليها الميل نفسه. وتدوير كل واجهة أمامية حول امتداد التطبّق بقدر انحدار التطبّق يعيد التطبّق، ومعه الواجهات الأمامية، إلى وضعيتها الأصلية.

\textbf{23.} تتدفّق التيارات المدّية داخلة وخارجة، وتبني مجموعات تنحدر في اتجاهين متعاكسين، كما في تطبّق متقاطع على شكل عظمة السمكة؛ وكان مخطط الوردة سيكون له فصّان متباعدان بنحو \ang{180}. والقراءات الاثنتا عشرة تكوّن فصًّا واحدًا، و$\bar R$ قريب من 1: تيار نهري يجري في اتجاه واحد.

\textbf{24.} $\bar\alpha = \operatorname{atan2}(1.698, 11.248) = \ang{8.6}$: متوسط \omterm{def:b1:clastic-sediments:cross-bedding}{التيار القديم} للسمت للوحدة~3 هو \ang{9}، شرق الشمال بقليل.
\end{solution}
